% !TeX root = ../../elegantbook-cn.tex
\begin{example}
（2026东城二模）已知函数 $f(x)=\ln x$，$g(x)=ax^2+b$．当 $m>0$ 时，曲线 $y=f(x)$ 在点 $(m,f(m))$ 处的切线为 $l_1$，曲线 $y=g(x)$ 在点 $(m,g(m))$ 处的切线为 $l_2$，且 $l_1\perp l_2$．\\
（1）求 $a$ 的值；\\
（2）当 $b>-\dfrac32$ 时，求证：$l_1$ 与 $l_2$ 的交点位于 $y$ 轴右侧；\\
（3）已知 $b>0$，$l_1$ 与 $y$ 轴交于点 $A$，$l_2$ 与 $y$ 轴交于点 $B$．若存在 $m\in(0,\mathrm{e})$，使得 $|OA|=\mathrm{e}|OB|$，求 $b$ 的最大值．
\end{example}

\begin{example}
（2025东城二模）设函数 $f(x)=(x^2+ax)\ln x$，其中 $a\in\mathbb{R}$．\\
（1）当 $a>0$ 时，求 $f(x)$ 的零点；\\
（2）当 $a=-1$ 时，证明：\\
（i）$1$ 为 $f(x)$ 的极小值点；\\
（ii）对于任意 $m\in\left(\dfrac12,1\right)$，存在 $n\in(1,2)$，使得曲线 $y=f(x)$ 在点 $A(m,f(m))$ 处的切线斜率与在点 $B(n,f(n))$ 处的切线斜率互为相反数．
\end{example}

\begin{example}
（2024东城二模）已知函数 $f(x)=x\sin2x+\cos2x$．\\
（1）求曲线 $y=f(x)$ 在点 $\left(-\dfrac\pi4,f\left(-\dfrac\pi4\right)\right)$ 处的切线方程；\\
（2）求函数 $f(x)$ 在区间 $\left[-\dfrac{2\pi}{3},\dfrac{5\pi}{6}\right]$ 上的极值点个数．
\end{example}

\begin{example}
（2023东城二模）已知函数 $f(x)=\mathrm{e}^x\sin x-2x$．\\
（1）求曲线 $y=f(x)$ 在点 $(0,f(0))$ 处的切线方程；\\
（2）求 $f(x)$ 在区间 $[-1,1]$ 上的最大值；\\
（3）设实数 $a$ 使得 $f(x)+x>a\mathrm{e}^x$ 对 $x\in\mathbb{R}$ 恒成立，写出 $a$ 的最大整数值，并说明理由．
\end{example}

\begin{example}
（2022东城二模）已知函数 $f(x)=x+\dfrac{2a^2}{x}+a\ln x$（$a\in\mathbb{R}$）．\\
（1）当 $a=1$ 时，求曲线 $y=f(x)$ 在点 $(1,f(1))$ 处的切线方程；\\
（2）当 $x\in[\mathrm{e},+\infty)$ 时，曲线 $y=f(x)$ 在 $x$ 轴的上方，求实数 $a$ 的取值范围．
\end{example}

\begin{example}
（2021东城二模）已知函数 $f(x)=\sin x-(x-a)\cos x$，其中 $a\in\left(-\dfrac\pi2,\dfrac\pi2\right)$．\\
（1）若曲线 $y=f(x)$ 在 $x=a$ 处的切线过点 $\left(0,\dfrac{\sqrt3}{2}\right)$，求 $a$ 的值；\\
（2）若 $f(x)>a^3-1$ 对 $x\in\left[-\dfrac\pi2,\dfrac\pi2\right]$ 恒成立，求 $a$ 的取值范围．
\end{example}

\begin{example}
（2026西城二模）已知函数 $f(x)=\dfrac{1+\ln x}{ax}$，其中 $a>0$．\\
（1）当 $a=1$ 时，求曲线 $y=f(x)$ 在点 $(1,f(1))$ 处的切线方程；\\
（2）对于 $x\in(0,+\infty)$，讨论 $f(x)$ 与 $f\left(\dfrac1x\right)$ 的大小；\\
（3）当 $0<a<1$ 时，证明：方程 $f(x)=1$ 存在两个根 $x_1,x_2$，且 $x_1x_2>1$．
\end{example}

\begin{example}
（2025西城二模）已知函数 $f(x)=a(x-1)\mathrm{e}^x-\ln x$，其中 $a>0$．\\
（1）若曲线 $y=f(x)$ 在点 $(1,f(1))$ 处的切线经过点 $(2,2)$，求 $a$ 的值；\\
（2）证明：函数 $f(x)$ 存在极小值；\\
（3）记函数 $f(x)$ 的最小值为 $g(a)$，求 $g(a)$ 的最大值．
\end{example}

\begin{example}
（2024西城二模）已知函数 $f(x)=4a\sin x+(a^2+1)x^2$，其中 $a\geqslant0$．\\
（1）若 $f(x)$ 在 $x=0$ 处取得极小值，求 $a$ 的值；\\
（2）当 $a=1$ 时，求 $f(x)$ 在区间 $\left[\dfrac\pi2,\pi\right]$ 上的最大值；\\
（3）证明：$f(x)$ 有且只有一个极值点．
\end{example}

\begin{example}
（2023西城二模）已知函数 $f(x)=x^2+\ln(x+1)$．\\
（1）求 $f(x)$ 在区间 $\left[-\dfrac12,0\right]$ 上的最大值和最小值；\\
（2）若 $(\mathrm{e}^x+a\cos x)f(x)\geqslant0$ 恒成立，求实数 $a$ 的值．
\end{example}

\begin{example}
（2022西城二模）已知函数 $f(x)=\dfrac{\ln x+a}{x+1}$．\\
（1）若 $f'(1)=\dfrac14$，求 $a$ 的值；\\
（2）当 $a>2$ 时，\\
①求证：$f(x)$ 有唯一的极值点 $x_1$；\\
②记 $f(x)$ 的零点为 $x_0$，是否存在 $a$ 使得 $\dfrac{x_1}{x_0}\leqslant \mathrm{e}$？说明理由．
\end{example}

\begin{example}
（2021西城二模）已知函数 $f(x)=\ln x+bx+c$，$g(x)=kx^2+2$，$f(x)$ 在 $x=1$ 处取得极大值 $1$．\\
（1）求 $b$ 和 $c$ 的值；\\
（2）当 $x\in[1,+\infty)$ 时，曲线 $y=f(x)$ 在曲线 $y=g(x)$ 的上方，求实数 $k$ 的取值范围；\\
（3）设 $k=1$，证明：存在两条与曲线 $y=f(x)$ 和 $y=g(x)$ 都相切的直线．
\end{example}

\begin{example}
（2026朝阳二模）已知函数 $f(x)=x-\dfrac1x-a\ln x$（$a\in\mathbb{R}$）．\\
（1）当 $a=1$ 时，求曲线 $y=f(x)$ 在点 $(1,f(1))$ 处的切线方程；\\
（2）当 $a=\mathrm{e}-\dfrac{1}{\mathrm{e}}$ 时，求满足 $f(x)>0$ 的 $x$ 的取值范围；\\
（3）当 $a\leqslant2$ 时，判断曲线 $y=f(x)$ 上是否存在两个不同的点关于点 $(1,0)$ 对称，并说明理由．
\end{example}

\begin{example}
（2025朝阳二模）已知函数 $f(x)=\dfrac{a\ln x+1}{\mathrm{e}^{x-1}}$（$a\geqslant0$）．\\
（1）若 $a=0$，求函数 $f(x)$ 在区间 $[1,+\infty)$ 上的最大值；\\
（2）若 $f(x)$ 在区间 $(0,1)$ 上存在单调递减区间，求 $a$ 的取值范围；\\
（3）若 $f(x)$ 存在极值点 $x_0$，且 $f(x_0)=1$，求 $a$ 的值．
\end{example}

\begin{example}
（2024朝阳二模）已知函数 $f(x)=ax-\ln(1-x)$（$a\in\mathbb{R}$）．\\
（1）求曲线 $y=f(x)$ 在点 $(0,f(0))$ 处的切线方程；\\
（2）若 $f(x)\geqslant0$ 恒成立，求 $a$ 的值；\\
（3）若 $f(x)$ 有两个不同的零点 $x_1,x_2$，且 $|x_2-x_1|>\mathrm{e}-1$，求 $a$ 的取值范围．
\end{example}

\begin{example}
（2023朝阳二模）已知函数 $f(x)=\ln x+2a\sqrt{x}$（$a\in\mathbb{R}$）．\\
（1）当 $a=1$ 时，\\
（i）求曲线 $y=f(x)$ 在点 $(1,f(1))$ 处的切线方程；\\
（ii）证明：$f(x)\leqslant2x$；\\
（2）若函数 $h(x)=f(x)-2x$ 的极大值大于 $0$，求 $a$ 的取值范围．
\end{example}

\begin{example}
（2022朝阳二模）已知函数 $f(x)=x\sin x+\cos x$．\\
（1）当 $x\in(0,\pi)$ 时，求函数 $f(x)$ 的单调区间；\\
（2）设函数 $g(x)=-x^2+2ax$．若对任意 $x_1\in[-\pi,\pi]$，存在 $x_2\in[0,1]$，使得 $\dfrac1{2\pi}f(x_1)\leqslant g(x_2)$ 成立，求实数 $a$ 的取值范围．
\end{example}

\begin{example}
（2021朝阳二模）已知函数 $f(x)=(x-1)\mathrm{e}^x-\dfrac12ax^2+1$，$a\in\mathbb{R}$．\\
（1）当 $a=0$ 时，求曲线 $y=f(x)$ 在点 $(1,f(1))$ 处的切线方程；\\
（2）判断函数 $f(x)$ 的极值点的个数，并说明理由；\\
（3）若对任意 $x\in\mathbb{R}$，$f(x)\geqslant0$ 恒成立，求 $a$ 的取值范围．
\end{example}

\begin{example}
（2026丰台二模）已知函数 $f(x)=(x-1)\ln x+ax^2$，$f'(x)$ 是 $f(x)$ 的导函数．\\
（1）若 $f'(1)-f(1)=1$，求 $a$ 的值；\\
（2）若 $f'(x)$ 存在最大值，求 $a$ 的取值范围；\\
（3）设 $f'(x)$ 在 $x=1$ 处取得最大值．直线 $l$ 是曲线 $y=f(x)$ 在点 $P(t,f(t))$（$t>1$）处的切线，且与 $y$ 轴交于点 $Q$，$O$ 为坐标原点．是否存在点 $P$，使 $\triangle OPQ$ 的面积为 $\dfrac14$？若存在，求 $t$ 的值；若不存在，说明理由．
\end{example}

\begin{example}
（2025丰台二模）已知函数 $f(x)=(x^3+2x^2)\mathrm{e}^{ax+b}$（$a,b\in\mathbb{R}$）在点 $(-1,f(-1))$ 处的切线方程为 $y=1$．\\
（1）求 $a,b$ 的值；\\
（2）求 $f(x)$ 的单调区间；\\
（3）若 $f(x)\leqslant-x^2-2x$，求 $x$ 的取值范围．
\end{example}

\begin{example}
（2024丰台二模）已知函数 $f(x)=a^2x+2a\sqrt{x}-2\ln x$（$a\neq0$）．\\
（1）当 $a=1$ 时，求曲线 $y=f(x)$ 在点 $(1,f(1))$ 处的切线方程；\\
（2）若函数 $f(x)$ 有两个零点，求 $a$ 的取值范围．
\end{example}

\begin{example}
（2023丰台二模）已知函数 $f(x)=(\sqrt{x}+a)\mathrm{e}^x$（$a\in\mathbb{R}$）．\\
（1）当 $a=0$ 时，求曲线 $y=f(x)$ 在点 $(1,f(1))$ 处的切线方程；\\
（2）若 $f(x)$ 是增函数，求 $a$ 的取值范围；\\
（3）证明：$f(x)$ 有最小值，且最小值小于 $f(1)$．
\end{example}

\begin{example}
（2022丰台二模）已知函数 $f(x)=\dfrac{ax+1}{\mathrm{e}^x}$．\\
（1）当 $a=1$ 时，求 $f(x)$ 的单调区间和极值；\\
（2）当 $a\geqslant1$ 时，求证：$f(x)\leqslant(a-1)x+1$；\\
（3）直接写出 $a$ 的一个取值范围，使得 $f(x)\geqslant ax^2+(a-1)x+1$ 恒成立．
\end{example}

\begin{example}
（2021丰台二模）已知函数 $f(x)=(x^2-2ax)\ln x-\dfrac12x^2+2ax$（$a\in\mathbb{R}$）．\\
（1）若 $a=0$，求 $f(x)$ 的最小值；\\
（2）求函数 $f(x)$ 的单调区间．
\end{example}

\begin{example}
（2026海淀二模）已知函数 $f(x)=\dfrac{x+a}{\sin x+2}$，其中 $a\geqslant0$．\\
（1）当 $a=0$ 时，求曲线 $y=f(x)$ 经过点 $(0,0)$ 的切线条数；\\
（2）当 $\dfrac\pi2<a\leqslant2$ 时，求证：对任意 $x\in\left[-\dfrac\pi2,\dfrac\pi2\right]$，$f(x)<\dfrac{2a}{3}$；\\
（3）若关于 $x$ 的方程 $f'(x)=1$ 在区间 $\left[-\dfrac\pi2,\dfrac\pi2\right]$ 上有且仅有 $1$ 个解，直接写出 $a$ 的最小值．
\end{example}

\begin{example}
（2025海淀二模）已知函数 $f(x)=\mathrm{e}^{-x}\ln(x+1)-ax$．\\
（1）求曲线 $y=f(x)$ 在点 $(0,f(0))$ 处的切线方程；\\
（2）求 $f(x)$ 在区间 $(-1,0)$ 上的极值点个数；\\
（3）若 $s,t\in(-1,+\infty)$ 且 $s+t\leqslant0$ 时，都有 $f(s)+f(t)\leqslant0$ 成立，直接写出 $a$ 的取值范围．
\end{example}

\begin{example}
（2024海淀二模）已知函数 $f(x)=\ln(x-a)+2\sqrt{3a-x}$（$a>0$）．\\
（1）当 $a=1$ 时，\\
①求曲线 $y=f(x)$ 在点 $(2,f(2))$ 处的切线方程；\\
②求证：函数 $f(x)$ 恰有一个零点；\\
（2）若 $f(x)\leqslant\ln a+2a$ 对 $x\in(a,3a)$ 恒成立，求 $a$ 的取值范围．
\end{example}

\begin{example}
（2023海淀二模）已知函数 $f(x)=\sqrt{x}\ln x$．\\
（1）求曲线 $y=f(x)$ 在点 $(1,f(1))$ 处的切线方程；\\
（2）求证：$f(x)<x$；\\
（3）若函数 $g(x)=f(x)+a(x^2-x)$ 在区间 $(1,+\infty)$ 上无零点，求 $a$ 的取值范围．
\end{example}

\begin{example}
（2022海淀二模）已知函数 $f(x)=\ln\dfrac{1-x}{2}+\dfrac ax$．\\
（1）当 $a=0$ 时，求曲线 $y=f(x)$ 在点 $(-1,f(-1))$ 处的切线方程；\\
（2）当 $a=-\dfrac12$ 时，求函数 $f(x)$ 的单调区间；\\
（3）当 $x<0$ 时，$f(x)\geqslant\dfrac12$ 恒成立，求 $a$ 的取值范围．
\end{example}

\begin{example}
（2021海淀二模）已知函数 $f(x)=x-a\ln x$．\\
（1）求曲线 $y=f(x)$ 在点 $(1,f(1))$ 处的切线方程；\\
（2）求 $f(x)$ 的单调区间；\\
（3）若关于 $x$ 的方程 $x-a\ln x=0$ 有两个不相等的实数根，记较小的实数根为 $x_0$，求证：$(a-1)x_0>a$．
\end{example}
